%%%%%%%%%%%%%%%%%%%%%%%%%%%%%%%%%%%%%%%%%%%%%%%%%%%%%%%%%%%%%%%%%%%%%%%%%
%% Linear Algebra -- course project report template, Autumn 2026
%% Ukrainian Catholic University, Faculty of Applied Sciences,
%% Data Science Master Programme
%%
%% The title page is adapted from a template by WikiBooks, distributed
%% by LaTeXTemplates.com under CC BY-NC-SA 3.0.
%%
%% HOW TO USE THIS TEMPLATE
%%   1. Fill in the block marked "FILL IN" below: the stage of the report,
%%      the title, the authors, the project card and the links.
%%   2. One file serves all three milestones. Keep working in the same
%%      file: interim report 1 grows into interim report 2, and that grows
%%      into the final report.
%%   3. Replace the grey guidance boxes with your own text. Before you
%%      submit, set \showguidefalse so that any box you forgot disappears.
%%   4. Do NOT change the font size, the paper size or the margins.
%%      Page limits are counted in this format.
%%   5. Compile with pdfLaTeX (the Overleaf default). The folder must
%%      contain UCU-Apps.png (the logo); the example figure uses an image
%%      shipped with every TeX distribution.
%%
%% PAGE LIMITS (counted automatically; see the end of the log file)
%%   Interim reports 1 and 2: 2 pages, project card included.
%%   Final report: at most 10 pages from the abstract to the end of the
%%   Conclusion, figures and tables included. The title page, the
%%   references and the appendices are not counted.
%%%%%%%%%%%%%%%%%%%%%%%%%%%%%%%%%%%%%%%%%%%%%%%%%%%%%%%%%%%%%%%%%%%%%%%%%

\documentclass[11pt,a4paper]{article}

%----------------------------------------------------------------------
%  Page layout: do not change
%----------------------------------------------------------------------
\usepackage[a4paper,left=20mm,right=20mm,top=20mm,bottom=25mm,
            headheight=14pt,marginparwidth=16mm]{geometry}

%----------------------------------------------------------------------
%  Packages
%----------------------------------------------------------------------
\usepackage[english]{babel}
% T1 encoding only together with the scalable Latin Modern fonts
% (standard on Overleaf); otherwise fall back to the default encoding.
\IfFileExists{lmodern.sty}{\usepackage[T1]{fontenc}\usepackage{lmodern}}{}
\usepackage{microtype}
\usepackage{amsmath,amssymb,amsthm}
\usepackage{graphicx}
\usepackage{booktabs,tabularx}
\usepackage{caption,subcaption}
\usepackage[dvipsnames]{xcolor}
\usepackage{fancyhdr}
\usepackage{refcount}
\usepackage{listings}
\usepackage[ruled,vlined,linesnumbered]{algorithm2e}
% Margin notes are useful while drafting. Add the option "disable"
% before you submit, or simply delete your \todo{} commands.
\usepackage[colorinlistoftodos,textsize=scriptsize]{todonotes}
\usepackage[colorlinks=true,linkcolor=MidnightBlue,citecolor=MidnightBlue,
            urlcolor=MidnightBlue]{hyperref}   % load hyperref last

%%%%%%%%%%%%%%%%%%%%%%%%%%%%%%%%%%%%%%%%%%%%%%%%%%%%%%%%%%%%%%%%%%%%%%%%%
%%  FILL IN
%%%%%%%%%%%%%%%%%%%%%%%%%%%%%%%%%%%%%%%%%%%%%%%%%%%%%%%%%%%%%%%%%%%%%%%%%

% Stage of the report: 1 = interim report 1 (scoping memo),
%                      2 = interim report 2 (working baseline),
%                      3 = final report
\newcommand{\stage}{1}

\newcommand{\reporttitle}{Low-Rank Spectral Fingerprints for AI-Generated Music Detection}
\newcommand{\shorttitle}{Low-rank spectrograms for AI-music detection}
\newcommand{\reportauthors}{Iryna Matvieieva \quad Petro Mozil}

% Project card: one line each. Keep it consistent with your registration;
% if something has changed, update it here and say so in the report.
\newcommand{\carddomain}{Audio forensics / music provenance}
\newcommand{\cardmethod}{STFT magnitude spectrograms; exact and randomized SVD}
\newcommand{\cardobject}{10-s mono clip; STFT magnitude $M\in\mathbb{R}_{\geq0}^{1025\times431}$; top $k=20$ singular values}
\newcommand{\cardalgorithms}{Exact NumPy/LAPACK SVD versus randomized SVD ($k=20$, $p=10$, one power step)}
\newcommand{\carddata}{Up to 500 paired clips from FakeMusicCaps and MusicCaps~\cite{FakeMusicCaps}}
\newcommand{\cardcode}{\url{https://github.com/pmozil/music-ai-spectra}}
\newcommand{\cardvideo}{\url{https://link-to-your-recording}} % final only

% Guidance boxes: set \showguidefalse before you submit.
\newif\ifshowguide
\showguidefalse

%%%%%%%%%%%%%%%%%%%%%%%%%%%%%%%%%%%%%%%%%%%%%%%%%%%%%%%%%%%%%%%%%%%%%%%%%
%%  Template machinery: no need to read further than the next block
%%%%%%%%%%%%%%%%%%%%%%%%%%%%%%%%%%%%%%%%%%%%%%%%%%%%%%%%%%%%%%%%%%%%%%%%%

% Notation used in the course. Add your own macros here.
\newcommand{\bR}{\mathbb{R}}
\newcommand{\bC}{\mathbb{C}}
\newcommand{\ba}{\mathbf{a}}
\newcommand{\bb}{\mathbf{b}}
\newcommand{\bu}{\mathbf{u}}
\newcommand{\bv}{\mathbf{v}}
\newcommand{\bx}{\mathbf{x}}
\newcommand{\by}{\mathbf{y}}
\newcommand{\bz}{\mathbf{z}}
\DeclareMathOperator{\rank}{rank}
\DeclareMathOperator{\cond}{cond}
\DeclareMathOperator{\diag}{diag}
\DeclareMathOperator{\trace}{trace}
\DeclareMathOperator{\nnz}{nnz}
\DeclareMathOperator*{\argmin}{argmin}

% Theorem-like environments, numbered within sections
\theoremstyle{plain}
\newtheorem{theorem}{Theorem}[section]
\newtheorem{proposition}[theorem]{Proposition}
\newtheorem{lemma}[theorem]{Lemma}
\theoremstyle{definition}
\newtheorem{definition}[theorem]{Definition}
\newtheorem{example}[theorem]{Example}
\theoremstyle{remark}
\newtheorem{remark}[theorem]{Remark}

% Code listings
\lstset{language=Python, basicstyle=\ttfamily\small, columns=fullflexible,
  keepspaces=true, frame=lines, numbers=left, numberstyle=\tiny\color{gray},
  commentstyle=\color{gray}, keywordstyle=\color{MidnightBlue}\bfseries,
  showstringspaces=false, breaklines=true}

% Guidance boxes
\newcommand{\guide}[1]{\ifshowguide\par\smallskip\noindent
  \colorbox{gray!12}{\parbox{\dimexpr\linewidth-2\fboxsep}{\small\itshape
  \textbf{Guidance.} #1}}\par\smallskip\fi}

% Stage names
\newcommand{\stagename}{\ifcase\stage\or Interim report 1: scoping memo\or
  Interim report 2: working baseline\or Final report\fi}

% Running heads
\pagestyle{fancy}
\fancyhf{}
\fancyhead[L]{\small\shorttitle}
\fancyhead[R]{\small\stagename}
\fancyfoot[C]{\small\thepage}
\renewcommand{\headrulewidth}{0.4pt}

% Project card
\newcommand{\projectcard}{%
  \noindent
  \begin{tabularx}{\linewidth}{@{}>{\raggedright\bfseries}p{3.0cm}
                                  >{\raggedright\arraybackslash}X@{}}
    \toprule
    Domain        & \carddomain \\
    Method family & \cardmethod \\
    LA object     & \cardobject \\
    Algorithms    & \cardalgorithms \\
    Data          & \carddata \\
    \ifnum\stage>1 Code & \cardcode \\ \fi
    \ifnum\stage=3 Recording & \cardvideo \\ \fi
    \bottomrule
  \end{tabularx}\par}

% Page counting: the main body runs from \bodystart to \bodyend
\newcommand{\pagelimit}{\ifnum\stage=3 10\else 2\fi}
\newcommand{\bodystart}{\label{body-start}}
\newcommand{\bodyend}{\label{body-end}}
\AtEndDocument{%
  \refused{body-start}\refused{body-end}%
  \edef\bodypages{\number\numexpr\getpagerefnumber{body-end}
    -\getpagerefnumber{body-start}+1\relax}%
  \typeout{}%
  \typeout{*** Main body: \bodypages\space page(s); limit \pagelimit.}%
  \typeout{}%
  \ifnum\bodypages>\pagelimit\relax
    \PackageWarningNoLine{LA-template}{Main body has \bodypages\space pages,
      above the limit of \pagelimit}%
  \fi}

%%%%%%%%%%%%%%%%%%%%%%%%%%%%%%%%%%%%%%%%%%%%%%%%%%%%%%%%%%%%%%%%%%%%%%%%%
\begin{document}

% Interim report 1: title and project card count toward the two-page limit.
\bodystart
\begin{center}
  {\Large\bfseries \reporttitle\par}
  \medskip
  {\large \stagename\par}
  \smallskip
  \begin{tabular}{c}\reportauthors\end{tabular}\par
  \smallskip
  {\small Linear Algebra, UCU Data Science Master Programme, \today\par}
\end{center}
\projectcard
\smallskip

\section{Problem and aim}
The detection of AI-generated music has proven a problem
for platforms, rights holders and listeners. In this work,
we intend to distinguish AI-generated audio from human-generated audio.
Given a short audio clip, we predict human-origin versus text-to-music-generated.
We do so by checking whether the low-rank structure of its STFT magnitude
spectrogram supports this distinction on held-out clips,
rather than create a universal detector. Prior work also warns that
strong controlled results may not generalize across manipulations
or datasets~\cite{Afchar,LiBench}, so we try to balance classes
and group each original/synthetic caption pair in the same data split.

We will sample up to 500 caption IDs from FakeMusicCaps and retrieve the
matching human-made MusicCaps clips. We will choose one generated clip per ID,
balance across the five generators where possible, and keep matched clips in
the same split. If original clips are inaccessible, we will report the smaller
paired sample rather than replace them with unmatched tracks. Success requires
held-out balanced accuracy $\geq0.65$ with a 95\% bootstrap confidence interval
whose lower bound exceeds $0.50$. The approximation must also satisfy
$\|p^{\rm rand}-p^{\rm exact}\|_2/\|p^{\rm exact}\|_2\leq10^{-2}$ on at least
95\% of clips; otherwise, randomized SVD will not be considered accurate
enough for this feature pipeline.

\section{Linear-algebra formulation}
Each ten-second clip is converted to mono at 22.05 kHz. A centred short-time
Fourier transform (STFT), using a 2048-sample window/FFT and hop 512, gives a
magnitude matrix $M\in\mathbb{R}_{\geq0}^{F\times T}$ with $F=1025$ and
$T\approx431$. Its singular value decomposition is

$$
  M=U\Sigma V^\top,\qquad \sigma_1\geq\sigma_2\geq\cdots\geq0.
$$

We classify the first 20 normalized squared singular values,
$p_i=\sigma_i^2/\|M\|_F^2$, using logistic regression. The linear-algebra task
is to compute or approximate this leading spectrum, which summarizes how
spectrogram energy is distributed across orthogonal components. It may discard
useful phase, local-frequency, and temporal cues, so classification is an
empirical feasibility test.

\section{Approaches and evaluation}
\textbf{Dense thin SVD (reference).} NumPy/LAPACK gives a numerically accurate
spectrum and is the reference for error measurement. For $F\geq T$, the cost
is $O(FT^2+T^3)$ flops and memory is $O(FT+T^2)$. Its advantage is reliability;
its drawback is that it computes the full spectrum, and these moderate matrix
sizes may leave little room for speedup.

\textbf{Randomized SVD.} Following Halko, Martinsson, and Tropp~\cite{Halko},
we will implement range finding with target rank $k=20$, oversampling $p=10$
($\ell=k+p=30$), QR orthogonalization, and one power step ($q=1$), then take
the SVD of the reduced matrix. Its cost is approximately
$O((2q+2)FT\ell+(F+T)\ell^2)$ flops and memory
$O(FT+(F+T)\ell)$. It focuses work on the leading subspace, but can lose
accuracy when singular values decay slowly or cluster; one power step may then
be insufficient, and extra steps increase runtime. We will implement this
randomized algorithm and the feature pipeline ourselves, while using NumPy /
LAPACK for the dense reference and scikit-learn for logistic regression.

We will compare feature error and wall-clock time, then report balanced accuracy,
ROC-AUC and bootstrap confidence intervals on a caption-grouped held-out set.
A bandwidth-only classifier is a diagnostic baseline: if it performs well,
labels may be explained by codec or spectral-bandwidth artifacts rather than
music structure. Frobenius normalization reduces sensitivity to overall gain,
but codec and bandwidth differences can still leak into labels; identical mono
conversion, resampling and loudness normalization will be applied to both classes.

\section{Risks and work plan}
The main risks are unavailable matched originals, a flat singular-value
spectrum, and leakage of generator or mastering artifacts into the classifier.
If the target accuracy is missed, we will report this negative result rather
than repeatedly tune on the test set. By 2 November, we will verify clip
availability and lock the preprocessing and split protocol. By 16 November
(interim report 2), we will implement the feature pipeline and randomized SVD,
validate its error and runtime against dense SVD on a subset, and run an initial
baseline. By 14 December, we will finish the paired evaluation (up to 500 IDs),
compare against the bandwidth baseline, calculate confidence intervals, and
discuss limitations in the final report. Iryna will lead data matching and
preprocessing; Petro will lead randomized SVD and numerical validation; both
will run the classification evaluation and write the report.

\bodyend

% The references are excluded from the two-page main-body limit.
\clearpage
\begin{thebibliography}{9}
\bibitem{FakeMusicCaps}
L. Comanducci, P. Bestagini and S. Tubaro,
``FakeMusicCaps: A Dataset for Detection and Attribution of Synthetic Music
Generated via Text-to-Music Models,'' \emph{Journal of Imaging}, vol. 11,
no. 7, article 242, 2025.
\href{https://doi.org/10.3390/jimaging11070242}{doi:10.3390/jimaging11070242}.

\bibitem{Afchar}
D. Afchar, G. Meseguer-Brocal and R. Hennequin,
``AI-Generated Music Detection and its Challenges,'' in
\emph{Proc. IEEE International Conference on Acoustics, Speech and Signal
Processing (ICASSP)}, 2025.
\href{https://doi.org/10.1109/ICASSP49660.2025.10890655}{doi:10.1109/ICASSP49660.2025.10890655}.

\bibitem{LiBench}
Y. Li, Q. Sun, H. Li, L. Specia and B. W. Schuller,
``Explainable detection of machine generated music and early systematic
evaluation,'' \emph{Scientific Reports}, vol. 16, article 13757, 2026.
\href{https://doi.org/10.1038/s41598-026-42133-7}{doi:10.1038/s41598-026-42133-7}.

\bibitem{Halko}
N. Halko, P. G. Martinsson and J. A. Tropp,
``Finding Structure with Randomness: Probabilistic Algorithms for Constructing
Approximate Matrix Decompositions,'' \emph{SIAM Review}, vol. 53, no. 2,
pp. 217--288, 2011.
\href{https://doi.org/10.1137/090771806}{doi:10.1137/090771806}.
\end{thebibliography}

\end{document}
